\subsection{Observables of the four-dimensional flow at positive flow time}
\label{rF13:sec}

The spatial flow of Proposition~\ref{prop:spatialflow} keeps reflection
positivity and breaks $\mathrm O(4)$ to $\mathrm O(3)\times\mathrm O(1)$; the
flow in all four directions~\cite{Lus10} keeps $B_4$ and mixes the two sides
of every reflection plane. This subsection treats the second one at a fixed
flow time $t>0$ in physical units, that is, at the lattice flow time
$s_k=t/a_k^2$. Its observables need no field renormalisation when they are
bounded: flowed Wilson loops of fixed physical size are bounded by $1$, so
their correlations have limits along subsequences for every tuning. What
the flow takes away is reflection positivity, and with it the bridge to the
transfer matrix and to Theorem~\ref{thm:avgtransfer}. We prove the
following, for every tuning $(a_k,\beta_k)$ and every flow time.
\begin{enumerate}[label=\textup{(\alph*)},leftmargin=*]
\item Every sequence of $k$ has a subsequence along which the joint laws of
the flowed packets of bounded local observables and of the flowed Wilson
loops of dyadic rectilinear loops converge; the limit is a random field with
bounded correlation kernels, invariant under $B_4$, under the translations
of $\R^4$ on the packets and under the dyadic translations on the loops
(Theorem~\ref{rF13:tight}). For the flowed action density
in physical units, $t^2E(t)$, the same holds under one bound on its mean,
(E$_t$), which is not proved.
\item Time reflection positivity of a flowed limit forces every continuous
smeared reflected two-point function to be a Laplace transform; a smeared two-point
function that is continuously differentiable across zero time separation
is then constant. So a flowed limit with the smoothness that the flow is
meant to produce is reflection positive only if its smeared two-point
functions do not decay in time (Theorem~\ref{rF13:smooth}); the flowed free
field is not reflection positive, and neither are its lattice flows for
large $k$ (Proposition~\ref{rF13:freeflow}).
\item Exact reflection positivity holds for observables flowed inside a
half-space with the links outside frozen. For them every subsequential limit
is a continuum limit in the sense of Definition~\ref{def:limit} without
renormalisation constants, and under the averaged condition~\ref{rE:I1} its
Hamiltonian has spectrum in $\{0\}\cup[m_*,\infty)$; these limits are states
prepared by a flow with a frozen boundary, not correlation functions of a
field (Theorem~\ref{rF13:half}). For the full flow
reflection positivity is proved only up to errors of order $2^{-d/a_k}$ at lattice
flow times $s_k\le d/(7\kappa a_k)$, that is at physical flow times tending
to zero (Proposition~\ref{rF13:approx}). Under a locality property of the
flow on the physical scale, (Loc$_t$), which is not proved, reflection
positivity holds asymptotically far from the plane and, under the averaged
condition, the full flowed limit clusters in time
(Proposition~\ref{rF13:cluster}).
\item $\mathrm O(4)$ invariance of a flowed limit holds if and only if one
rotation of infinite order is an asymptotic symmetry, and under (E$_t$)
flowed Wilson loops depend continuously on the loop. The lattice structure
does not give the rotation (Remark~\ref{rF13:o4nogo}).
\item An area law for flowed Wilson loops in physical units excludes every
limit whose flowed Wilson loops are those of an abelian Gaussian gauge field
with finite line integrals of its covariance, in particular the flat limit
and the heat-flowed free Maxwell field; the heat-flowed Proca field shows
that exponential decay does not give the area law
(Theorem~\ref{rF13:area}).
\end{enumerate}
The flow-time-zero limit $t\to0$, which would connect these limits to the
Osterwalder--Schrader framework of Sec.~\ref{ssec:limitE} and to
Theorem~\ref{thm:avgtransfer}, is the content of (E) and is not treated;
Sec.~\ref{rF13:zero} states what remains.

\subsubsection{The lattice flow}

Let $G\subset\U(N)$ be a closed subgroup with Lie algebra $\g\subset
\mathfrak u(N)$, and let $P_\g$ be the orthogonal projection of the
$N\times N$ complex matrices onto $\g$ for the real inner product
$\operatorname{Re}\Tr(A^*B)$; on $\g$ this is the bi-invariant metric
$\ip XY=-\Tr(XY)$. Write $E$ for the links of $\Z^4$ with the positive
orientation, $V_{-e}=V_e^{-1}$, and for a plaquette $p$ containing $e$ let
$M_{p,e}(V)$ be the product of the four links of $p$ in the order of a
traversal that starts with $e$ in its positive direction. For a link $e$ put
\begin{equation}\label{rF13:eq:Y}
Y_e(V)=-\tfrac12\sum_{p\ni e}P_\g\bigl(M_{p,e}(V)-M_{p,e}(V)^*\bigr),
\end{equation}
the sum over the six plaquettes containing $e$. The \emph{lattice flow}
$\Phi_s(U)=V(s)$ is the solution of
\begin{equation}\label{rF13:eq:flow}
\dot V_e=Y_e(V)\,V_e\quad(e\in E),\qquad V(0)=U .
\end{equation}
The function $N-\ReTr V_p$ has gradient $\frac12P_\g(M_{p,e}-M_{p,e}^*)V_e$ in
the variable $V_e$ for the metric $\ip XY$ transported by right
multiplication, since its derivative along $V_e\mapsto e^{\varepsilon X}V_e$ is
$-\operatorname{Re}\Tr(XM_{p,e})=-\operatorname{Re}\Tr\bigl((M_{p,e}^*)^*X\bigr)$
and the Hermitian part of $M_{p,e}^*$ is orthogonal to $\g$. So
\eqref{rF13:eq:flow} is the gradient flow of the Wilson action
$S(V)=\sum_p(N-\ReTr V_p)$, the sum over unoriented plaquettes. For
$G=\SU(N)$ it is the Wilson flow (1.4) of~\cite{Lus10} in lattice units:
there the generators satisfy $\Tr(T^aT^b)=-\frac12\delta^{ab}$, so
$\partial_{x,\mu}=T^a\partial^a_{x,\mu}$ is the gradient for the metric
$-2\Tr(XY)$, and the action of~\cite{Lus10} sums over oriented plaquettes,
which doubles $S$. The flow time $s$ of~\eqref{rF13:eq:flow} is
dimensionless; at spacing $a$ it corresponds to the physical flow time
$t=a^2s$ of the continuum flow $\dot B_\mu=D_\nu G_{\nu\mu}$ of~\cite[(1.1)]{Lus10}.
Narayanan and Neuberger smoothed lattice Wilson loops by iterated APE
smearing with a number of steps proportional to a physical length squared,
identified the continuum version of this smearing with the gradient flow of
the Yang--Mills action, the Langevin equation without noise~\cite[Sec.~4.1,
(4.1)]{NN06}, and found
numerically that the smoothed loops have a finite and nontrivial continuum
limit~\cite{NN06}; L\"uscher showed that the expectation of the flowed action
density needs no renormalisation at one loop and gave numerical evidence
for its continuum limit~\cite{Lus10}, and L\"uscher and Weisz extended the
finiteness to all orders of perturbation theory~\cite{LW11}. None of this is
used below. What is used of the flow is collected in the next lemma, which
we prove; on a finite lattice its first part is~\cite[Sec.~3.1]{Lus09}.

Two links are \emph{neighbours} if they lie in a common plaquette; each link
has $18$ neighbours, and $\dist(e,F)$ denotes the graph distance from a link
$e$ to a set $F$ of links. For $h\in\Z$ let $E_h$ be the set of links with
both endpoints at times $\ge h$, and let $\Phi^{(h)}_s$ be the
\emph{half-space flow}: the links of $E_h$ solve~\eqref{rF13:eq:flow} with
$Y_e$ replaced by the sum over the plaquettes contained in $E_h$, and all
other links are held fixed. Translations and elements of $B_4$ act on
configurations as in the proof of Lemma~\ref{rE:torus}(ii), carrying each
link variable to the variable of the image link or to its inverse, and a
gauge transformation $g$ acts by $(gV)_{(x,y)}=g_xV_{(x,y)}g_y^{-1}$. Put
$e_x(V)=\sum_{\mu<\nu}(N-\ReTr V_{p(x;\mu,\nu)})$, the action of the six
plaquettes with lower corner $x$.

\begin{lemma}\label{rF13:flow}
There is a constant $\kappa$ depending only on $G$ such that the following
hold, with $\|\cdot\|$ the operator norm.

\textup{(i)} For every configuration $U\in G^E$ the system~\eqref{rF13:eq:flow}
has a unique solution for $s\in\R$, with values in $G^E$; $\Phi_s\circ
\Phi_{s'}=\Phi_{s+s'}$. The same holds for $\Phi^{(h)}_s$.

\textup{(ii)} $\Phi_s$ commutes with gauge transformations, translations and
$B_4$; $\Phi^{(h)}_s$ commutes with gauge transformations, with the
translations and elements of $B_4$ that fix the time coordinate, and
$\tau_n\Phi^{(h)}_s\tau_{-n}=\Phi^{(h+n)}_s$ for the time translation $\tau_n$.

\textup{(iii)} If $U,U'$ agree on every link at distance $<R$ from $e$, then
$\|\Phi_s(U)_e-\Phi_s(U')_e\|\le2e^{\kappa s}(\kappa s)^R/R!$ for $s\ge0$.

\textup{(iv)} If both endpoints of $e$ are at times $\ge h+R$, $R\ge1$,
then $\|\Phi_s(U)_e-\Phi^{(h)}_s(U)_e\|\le2e^{\kappa s}(\kappa s)^R/R!$
for $s\ge0$ and every $U$. If $\kappa s\le R/7$, the right sides of
\textup{(iii)} and \textup{(iv)} are at most $2^{1-R}$.

\textup{(v)} Let $O$ be a continuous function of the links of a finite set.
Then $O\circ\Phi_s$ and $O\circ\Phi^{(h)}_s$ are uniform limits of
continuous functions of finitely many links, gauge invariant if $O$ is, and
$O\circ\Phi^{(h)}_s$ depends only on the links of $E_h$.

\textup{(vi)} If $\omega$ is a probability measure on $G^E$ invariant under
translations, then $s\mapsto\omega(e_x\circ\Phi_s)$ is differentiable and
nonincreasing, with $\frac d{ds}\omega(e_x\circ\Phi_s)=-\sum_\mu\omega\bigl(
|Y_{(x,\mu)}|^2\circ\Phi_s\bigr)$, $|Y|^2=-\Tr(Y^2)$.
\end{lemma}

\begin{proof}
(i) Let $\mathcal B$ be the Banach space of bounded maps from $E$ to the
complex $N\times N$ matrices with the norm $\sup_e\|V_e\|$. By
\eqref{rF13:eq:Y}, $V\mapsto(Y_e(V)V_e)_e$ is a map from $\mathcal B$ to
$\mathcal B$ whose $e$-component is a polynomial of degree five in the
entries of the links of the six plaquettes at $e$ and their conjugates,
with coefficients independent of $e$; on the ball $\{\sup_e\|V_e\|\le2\}$
it is bounded and Lipschitz with constants independent of $e$, because
every component involves $19$ links. The Picard--Lindel\"of theorem in
$\mathcal B$ gives a unique local solution. For every $V\in\mathcal B$ the
matrix $Y_e(V)$ lies in $\g$ and is anti-Hermitian, since $M-M^*$ is
anti-Hermitian and $P_\g$ maps into $\g$; so $\frac d{ds}(V_e^*V_e)=V_e^*
(Y_e^*+Y_e)V_e=0$, the solution stays unitary, and $V_e(s)=W_e(s)U_e$ where
$W_e$ solves $\dot W_e=Y_e(V(s))W_e$, $W_e(0)=1$, a curve in $G$ because its
velocity is right-translated from $\g$. A solution with values in the
unitary matrices stays in the ball, so it extends to all $s$; uniqueness
gives the group law. For $\Phi^{(h)}_s$ the same argument applies to the
links of $E_h$, with the frozen links as parameters.

(ii) A gauge transformation maps $M_{p,e}$ to $g_xM_{p,e}g_x^{-1}$, $x$ the
initial point of $e$, and $P_\g$ commutes with $\mathrm{Ad}(g_x)$, so
$Y_e(gV)=g_xY_e(V)g_x^{-1}$ and $t\mapsto g\Phi_s(U)$ solves
\eqref{rF13:eq:flow} with initial value $gU$. A translation or an element
of $B_4$ permutes the plaquettes and maps $M_{p,e}$ to the corresponding
product, or to the inverse of a conjugate of it when the orientation of $e$
is reversed; $P_\g$ commutes with $A\mapsto-A^*$, and inversion is an
isometry of the bi-invariant metric, so the vector field is equivariant.
Uniqueness gives the claims; for $\Phi^{(h)}_s$ the maps must preserve
$E_h$, or move it to $E_{h+n}$.

(iii), (iv) Both are instances of one estimate. Let $V,V'$ be solutions of
two systems of the form~\eqref{rF13:eq:flow} with values in $G^E$, let $I$
be a set of links on which the two systems have the same equations and
$V(0)=V'(0)$, and put $d_{e'}(s)=\|V_{e'}(s)-V'_{e'}(s)\|$, so $d\le2$.
Expanding the products $M_{p,e'}$ link by link, using that unitary
matrices have norm one and that $P_\g$ is bounded, one finds $c$ depending
only on $G$ with $\|Y_{e'}(V)V_{e'}-Y_{e'}(V')V'_{e'}\|\le c\sum_{e''}
d_{e''}$, the sum over $e'$ and its $18$ neighbours; put $\kappa=19c$.
Then for $e'\in I$
\[
d_{e'}(s)\le\int_0^s\Bigl(c\,d_{e'}+c\sum_{e''\sim e'}d_{e''}\Bigr)dr .
\]
Write this as $u\le\int_0^s(Mu+Mb)$, where $u$ is the restriction of $d$ to
$I$, $b\le2$ its restriction to the complement $B$, and $M\ge0$ is the
matrix with row sums at most $\kappa$ that has entry $c$ at $(e',e'')$ when
$e''=e'$ or $e''\sim e'$. Iterating $n$ times,
\[
u_e(s)\le\sum_{j=1}^n\frac{s^j}{j!}\bigl(M_{II}^{j-1}M_{IB}\,2\bigr)_e
+\frac{(\kappa s)^n}{n!}\,2 ,
\]
and the $j$-th term vanishes unless there is a path of neighbours of length
$j$ from $e$ into $B$, that is unless $j\ge\dist(e,B)$; it is at most
$2(\kappa s)^j/j!$. Letting $n\to\infty$, $d_e(s)\le2\sum_{j\ge\dist(e,B)}
(\kappa s)^j/j!$, and $\sum_{j\ge R}x^j/j!\le e^xx^R/R!$. For (iii) take
$I$ the links at distance $<R$ from $e$. For (iv) take $I$ the links with
both endpoints at times $\ge h+1$: every plaquette containing such a link
lies in $E_h$, so the equations of the two flows agree there; a step to a
neighbour lowers the smaller time of the endpoints by at most one, so
$\dist(e,B)\ge R$. Finally $R!\ge(R/e)^R$ gives
$e^xx^R/R!\le(e^{8/7}/7)^R<2^{-R}$ for $x\le R/7$.

(v) Let $\Phi^R_s$ be the flow of the links at distance $<R$ from the
support of $O$, all others fixed. It is a finite system, its solution is a
continuous function of the finitely many links that enter it, and it
commutes with gauge transformations by the argument of (ii). The estimate
of (iii), with $I$ the links at distance $<R$ from the support and the two
systems $\Phi_s$ and $\Phi^R_s$, gives $\sup_U\|\Phi_s(U)_e-\Phi^R_s(U)_e\|
\to0$ for every $e$ in the support, and $O$ is uniformly continuous. The same
applies to $\Phi^{(h)}_s$, which never involves links outside $E_h$.

(vi) The function $e_x\circ\Phi_s$ is continuously differentiable in $s$
with derivative bounded uniformly in $U$, so the derivative may be taken
under $\omega$. The derivative of $N-\ReTr V_p$ along the flow is
$\sum_{e\in p}\ip{\nabla_e(N-\ReTr V_p)}{Y_eV_e}$, and $Y_eV_e=-\nabla_eS$.
By translation invariance of $\omega$ and (ii), $\omega(F\circ\Phi_s)$ is
unchanged when $F$ is translated, so the sum over the plaquettes $p$ with
lower corner $x$ and over $e\in p$ may be rearranged into the sum over the
four links $e$ starting at $x$ and the plaquettes $p\ni e$. Since
$\sum_{p\ni e}\nabla_e(N-\ReTr V_p)=\nabla_eS$, the derivative is
$-\sum_\mu\omega(|\nabla_{(x,\mu)}S|^2\circ\Phi_s)$, and $|\nabla_eS|=|Y_e|$.
\end{proof}

Part (iv) is the only locality of the flow that holds for every
configuration, and it is ballistic: the influence of a link reaches distance
$R$ at flow time of order $R$. At $s=t/a_k^2$ it covers $t/a_k^2$ lattice
spacings, the physical distance $t/a_k$, which diverges. The linearised flow
acts diffusively, over the physical distance $\sqrt{8t}$, but no
bound of that form holds for all configurations by the argument above: every
critical point of $S$ is fixed by the flow, and at a critical point whose
Hessian has a negative eigenvalue the linearised flow amplifies the
corresponding mode by a factor exponential in $s$. A locality on the scale
$\sqrt t$ uniform in $k$ is therefore a statement about typical
configurations under $\omega_k$; it is the input (Loc$_t$) below and is not
proved.
\subsubsection{Tightness and subsequential limits}

Throughout, $a_k=2^{-k}a_0$, the states $\omega_k$ are probability measures on
$G^E$ invariant under the translations of $\Z^4$ and under $B_4$ (the torus
limits of Lemma~\ref{rE:torus} are examples, at any couplings $\beta_k$), and
$s_k\ge0$ are lattice flow times; the case of interest is $s_k=t/a_k^2$ with
$t>0$ fixed, but nothing in this paragraph uses it. Let $\mathcal O$ be a
countable set of bounded continuous gauge-invariant functions of the links
of finite sets, closed under complex conjugation and under $O\mapsto O^g$,
$g\in B_4$, where $O^g(V)=O(g^{-1}V)$; write $O_x$ for the translate of $O$
to $x\in\Z^4$, so that $O^g_{gx}(gV)=O_x(V)$. The flowed packets are
\[
O_k(f)=\sum_{x\in\Z^4}a_k^4f(a_kx)\,O_x\circ\Phi_{s_k},
\]
for $f\in C_c(\R^4)$ or $f$ the indicator $\mathbf 1_Q$ of a half-open dyadic
cube $Q=a_02^{-m}(z+[0,1)^4)$, $z\in\Z^4$, of level $m$. Let $\mathcal C$ be
the set of closed rectilinear polygons in $\R^4$ with edges parallel to the
axes and vertices in $a_02^{-m}\Z^4$ for some $m$; for $k\ge m$ such a loop $C$
is a closed path of the lattice $a_k\Z^4$, and with $V=\Phi_{s_k}(U)$ and
$V(\gamma)$ the ordered product of the links along a lattice path $\gamma$,
\[
W_k(C)=N^{-1}\Tr V(C/a_k),\qquad|W_k(C)|\le1 ,
\]
the trace not depending on the starting point. The flowed action density in
physical units is
\[
\mathfrak E_k(x)=2s_k^2\,e_x\circ\Phi_{s_k},
\]
which at $s_k=t/a_k^2$ is $t^2E(t,a_kx)$ for the density $E=2\sum_{p\in
P_x}\ReTr\{1-V(p)\}$ of~\cite[(3.1)]{Lus10}, converted to physical units by
$a_k^{-4}$. It is bounded only by $24Ns_k^2$, and its packets
$\mathfrak E_k(f)=\sum_xa_k^4f(a_kx)\mathfrak E_k(x)$ need a bound:
\begin{enumerate}[label=\textup{(E$_t$)},leftmargin=*]
\item \emph{Bounded flowed action.} $K_E=\sup_k\omega_k(\mathfrak E_k(0))
<\infty$.
\end{enumerate}
In perturbation theory $\omega(t^2E(t))=3(N^2-1)g^2/(128\pi^2)+O(g^4)$ with
the coupling $g$ at the scale $(8t)^{-1/2}$~\cite[(2.32)]{Lus10}, so (E$_t$)
is the statement that this number stays bounded along the tuning. It is not
proved for any tuning with $\beta_k\to\infty$; Lemma~\ref{rF13:flow}(vi)
gives only $\omega_k(\mathfrak E_k(0))\le2s_k^2\omega_k(e_0)\le24Ns_k^2$.

\begin{theorem}\label{rF13:tight}
Every sequence of $k$ has a subsequence along which the following hold, on
a probability space that carries random variables $X^O(f)$, $O\in\mathcal
O$, $f\in C_c(\R^4)$, and $W(C)$, $C\in\mathcal C$.

\textup{(i)} For finitely many $O_j\in\mathcal O$, $f_j\in C_c(\R^4)$ and
$C_i\in\mathcal C$, the joint law of $(O_{j,k}(f_j),W_k(C_i))$ under
$\omega_k$ converges to that of $(X^{O_j}(f_j),W(C_i))$; in particular every
mixed moment converges.

\textup{(ii)} $f\mapsto X^O(f)$ is linear, $|X^O(f)|\le\|O\|_\infty\|f\|_{L^1}$,
$X^{\bar O}(\bar f)=\overline{X^O(f)}$ and $|W(C)|\le1$.

\textup{(iii)} For $O_1,\dots,O_n\in\mathcal O$ and $C_1,\dots,C_m\in\mathcal
C$ there is a kernel $S\in L^\infty((\R^4)^n)$ with $\|S\|_\infty\le\prod_j
\|O_j\|_\infty$ such that
\[
\E\Bigl[\prod_jX^{O_j}(f_j)\prod_iW(C_i)\Bigr]=\int S(x_1,\dots,x_n)
\prod_jf_j(x_j)\,dx .
\]

\textup{(iv)} For every dyadic $v\in\bigcup_ma_02^{-m}\Z^4$ the family
$(X^O(f(\cdot-v)),W(C+v))$ has the law of $(X^O(f),W(C))$; without loops
this holds for every $v\in\R^4$, and the kernels of \textup{(iii)} with $m=0$
are invariant under all translations of $\R^4$. For $g\in B_4$ the family
$(X^{O^g}(f\circ g^{-1}),W(gC))$ has the law of $(X^O(f),W(C))$.

\textup{(v)} The limit is a random field: the moments of \textup{(i)} are
those of a probability measure, so that, for instance,
$\sum_{i,j}\bar c_ic_j\E[\overline{X^{O_i}(f_i)}X^{O_j}(f_j)]\ge0$.

\textup{(vi)} Under \textup{(E$_t$)} the subsequence can be chosen so that
in addition the packets $\mathfrak E_k(f)$ converge jointly in law with the
families of \textup{(i)} to random variables $X^{\mathfrak E}(f)$; almost
surely $f\mapsto X^{\mathfrak E}(f)$ is integration against a positive Radon
measure, $\E X^{\mathfrak E}(f)\le K_E\int f$ for $f\ge0$, and the laws are
invariant as in \textup{(iv)}. If moreover $\sup_k\omega_k(\mathfrak E_k(0)^q)
<\infty$ for some $q>1$, the mixed moments of order less than $q$ in the
$\mathfrak E$-packets converge.
\end{theorem}

\begin{proof}
\emph{Cubes.} For $k\ge m$ the points $a_kx$ in a cube $Q$ of level $m$
number exactly $|Q|a_k^{-4}$, so $|O_k(\mathbf 1_Q)|\le\|O\|_\infty|Q|$, and
$O_k(\mathbf 1_Q)$ is the sum of $O_k(\mathbf 1_{Q'})$ over the $16$ cubes of
level $m+1$ in $Q$. The family of all $O_k(\mathbf 1_Q)$ and $W_k(C)$, with
$O\in\mathcal O$, dyadic $Q$, $C\in\mathcal C$ (set to $0$ for the finitely
many $k$ below the level), takes values in a countable product $K$ of closed
discs, a compact metrisable space. By Lemma~\ref{lem:helly} a subsequence
of the laws converges weakly to a probability measure $P$ on $K$; let
$X^O(Q)$ and $W(C)$ be the coordinates. The continuous function
$X^O(Q)-\sum_{Q'}X^O(Q')$ vanishes on the lattice laws for large $k$, so it
vanishes $P$-almost surely, and $|X^O(Q)|\le\|O\|_\infty|Q|$ almost surely,
the set being closed. On the event, of full measure, where these countably
many relations hold, $Q\mapsto X^O(Q)$ is finitely additive on dyadic cubes.

\emph{Continuous test functions.} For $f\in C_c(\R^4)$ let $\varpi_f$ be its
modulus of continuity, $\mathcal Q_m(f)$ the cubes of level $m$ that meet the
support of $f$, $c_Q$ the centre of $Q$, $\delta_m=2^{1-m}a_0$ their
diameter and $|\mathcal Q_m(f)|$ their total volume, bounded in $m$. For
$k\ge m$,
\[
\Bigl|O_k(f)-\sum_{Q\in\mathcal Q_m(f)}f(c_Q)O_k(\mathbf 1_Q)\Bigr|\le
\|O\|_\infty\varpi_f(\delta_m)|\mathcal Q_m(f)| ,
\]
and by additivity the same bound holds for the differences of the sums
\[\sum_{Q\in\mathcal Q_m(f)}f(c_Q)X^O(Q)\]
at levels $m\le m'$. So these sums
converge almost surely; call the limit $X^O(f)$. For a bounded Lipschitz
function $\varphi$ of finitely many variables, $\omega_k(\varphi(O_{j,k}(f_j),
W_k(C_i)))$ differs from its value at the level-$m$ sums by at most
$\mathrm{Lip}(\varphi)\sum_j\|O_j\|_\infty\varpi_{f_j}(\delta_m)|\mathcal
Q_m(f_j)|$ uniformly in $k\ge m$, the level-$m$ expressions converge by the
weak convergence of $P_k$, and the same bound holds in the limit; bounded
Lipschitz functions determine weak convergence, which gives (i). Linearity,
the bound by $\|f\|_{L^1}$ and the conjugation identity in (ii) pass from the
sums.

(iii) Fix the $O_j$ and the $C_i$. Let $F=\sum_Bc_B\mathbf 1_B$ be a function
on $(\R^4)^n$ that is a finite combination of products $B=Q_1\times\dots
\times Q_n$ of cubes of one level $m$, and put
\[
\Lambda(F)=\sum_Bc_B\,\E\Bigl[\prod_jX^{O_j}(Q_j)\prod_iW(C_i)\Bigr].
\]
By additivity $\Lambda$ does not depend on the level used to write $F$, and
$|\Lambda(F)|\le\prod_j\|O_j\|_\infty\sum_B|c_B||B|=\prod_j\|O_j\|_\infty
\|F\|_{L^1}$. Such $F$ are dense in $L^1((\R^4)^n)$, so $\Lambda$ extends to
$L^1$ with the same norm, and the duality $L^1{}^*=L^\infty$ gives $S$. For
$f_j\in C_c$ let $f^m_j=\sum_{Q\in\mathcal Q_m(f_j)}f_j(c_Q)\mathbf 1_Q$. Then
$\prod_jf^m_j(x_j)\to\prod_jf_j(x_j)$ in $L^1$, $X^{O_j}(f^m_j)\to X^{O_j}(f_j)$
with a deterministic bound, and dominated convergence gives the formula.

(iv) Let $v\in a_02^{-m}\Z^4$ and $k\ge m$, so $v/a_k\in\Z^4$. With
$\tau=\tau_{v/a_k}$, Lemma~\ref{rF13:flow}(ii) gives
$O_k(f(\cdot-v))(\tau U)=O_k(f)(U)$ and $W_k(C+v)(\tau U)=W_k(C)(U)$, and
$\omega_k$ is invariant under $\tau$; so the lattice laws agree and so do
the limits. Similarly, for $g\in B_4$ the substitution $y=gx$ gives
$O^g_k(f\circ g^{-1})(gU)=O_k(f)(U)$, and $W_k(gC)(gU)=W_k(C)(U)$ by the
definition of the action on oriented links. Without loops, the joint
characteristic function of finitely many $X^{O_j}(f_j(\cdot-v))$ is
continuous in $v\in\R^4$, because $\|f(\cdot-v)-f(\cdot-v')\|_{L^1}\to0$ as
$v'\to v$ and (ii) holds, and it is constant on the dense set of dyadic $v$.
For the kernels, $v\mapsto\int S(x+v)F(x)dx$ is continuous for $F\in L^1$
and constant on dyadic $v$.

(v) is a restatement of (i): the limits are expectations under $P$.

(vi) Add the coordinates $\mathfrak E_k(\mathbf 1_Q)\ge0$. By translation
invariance $\omega_k(\mathfrak E_k(\mathbf 1_Q))=|Q|\omega_k(\mathfrak E_k(0))
\le K_E|Q|$, so each coordinate is tight on $[0,\infty)$ by Markov's
inequality, and a countable family of tight coordinates is tight in the
product topology (choose compact sets with failure probabilities
$\varepsilon2^{-i}$). Prokhorov's theorem on the Polish space $[0,\infty)^{\mathbb
N}\times K$~\cite{Kal21} gives a subsequence. Additivity and positivity hold
almost surely as before, and $\E X^{\mathfrak E}(Q)\le\liminf_k\omega_k(
\mathfrak E_k(\mathbf 1_Q))\le K_E|Q|$, testing against $\min(x,M)$ and
letting $M\to\infty$. For $f\in C_c$ the error of the level-$m$ sums is now
bounded in $L^1(\omega_k)$ by $\varpi_f(\delta_m)K_E|\mathcal Q_m(f)|$, and
almost surely by $\varpi_f(\delta_m)X^{\mathfrak E}(\mathcal Q_m)$ in the limit,
which suffices for convergence in law by the same Lipschitz argument. On the
event where additivity and positivity hold for all cubes, $f\mapsto
X^{\mathfrak E}(f)$ is a positive linear functional on $C_c(\R^4)$, hence a
positive Radon measure. Invariance under dyadic translations is proved as in (iv). For $g\in B_4$, $e_x(gU)$ is the action of the six plaquettes at
the corner $g^{-1}x$ of the configuration $U$, which have lower corners
within lattice distance $4$ of $g^{-1}x$; so $\mathfrak E_k(f\circ g^{-1})
(gU)-\mathfrak E_k(f)(U)$ is a sum over plaquettes of $2s_k^2(N-\ReTr V_p)$
times differences of values of $f$ at points at distance at most $4a_k$,
bounded by $\varpi_f(4a_k)\sum_xa_k^4\mathbf 1_{\{\dist(a_kx,\operatorname{supp}f)\le4a_k\}}
\mathfrak E_k(x)$, whose expectation tends to zero; the laws therefore
have the same limits. Finally, for
$f\ge0$ Jensen's inequality for the probability weights $a_k^4f(a_kx)/
\sum_ya_k^4f(a_ky)$ gives $\mathfrak E_k(f)^q\le\|f\|_{1,k}^{q-1}\sum_xa_k^4
f(a_kx)\mathfrak E_k(x)^q$, so $\omega_k(\mathfrak E_k(f)^q)$ is bounded,
the powers of order less than $q$ are uniformly integrable, and their
expectations converge with the laws.
\end{proof}

The flow enters the proof only through Lemma~\ref{rF13:flow}(i),(ii),(v):
no renormalisation, no bound on the flow and no property of the coupling
is used for (i)--(v). The content of a limit therefore depends on what the
observables see, and here the bound (E$_t$) has a sharp consequence:
packets of observables of lattice size become deterministic, so that only
loops of physical size and the action density in physical units remain.

\begin{lemma}\label{rF13:stokes}
\textup{(i)} For $U\in\U(N)$, $\|1-U\|^2\le2(N-\ReTr U)$.

\textup{(ii)} If a based lattice path $\gamma'$ arises from $\gamma$ by
inserting at one of its vertices the boundary loop of a plaquette $p$
having that vertex as a corner, and cancelling backtracks, then
$\|V(\gamma')-V(\gamma)\|=\|1-V_p\|$. Consequently, if a closed path $\gamma$
in a coordinate plane is the boundary of a union $Y$ of plaquettes forming a
Young diagram, based at its inner corner, then $\|V(\gamma)-1\|\le\sum_{p\in
Y}\|1-V_p\|$.

\textup{(iii)} Under \textup{(E$_t$)} with $s_k\to\infty$, for every bounded
continuous gauge-invariant $O$ of the links of a finite set,
$\omega_k(|O\circ\Phi_{s_k}-O(\mathbf 1)|)\to0$, $\mathbf 1$ the configuration
with all links $1$; hence $X^O(f)=O(\mathbf 1)\int f$ almost surely.
\end{lemma}

\begin{proof}
(i) If $U$ has eigenvalues $e^{i\theta_j}$, then $\|1-U\|^2=\max_j
2(1-\cos\theta_j)\le\sum_j2(1-\cos\theta_j)$.

(ii) The holonomy of the inserted loop is $AV_p^{\pm1}A^{-1}$ with $A$
unitary, cancelling backtracks does not change holonomies, and
$V(\gamma')-V(\gamma)=V(\alpha)(AV_p^{\pm1}A^{-1}-1)V(\beta)$ for the
decomposition $\gamma=\alpha\beta$ at the vertex; unitary factors do not
change the norm. Removing an outer corner plaquette from a Young diagram
leaves a Young diagram, and the two boundaries differ by one such insertion
at a corner of the removed plaquette; induction on $|Y|$ gives the bound.

(iii) Let the links of $O$ lie in a box $\{0,\dots,L\}^4$. Choose the comb
tree of the box (all links of direction $0$; those of direction $1$ at time
$0$; those of direction $2$ with $y_0=y_1=0$; those of direction $3$ with
$y_0=y_1=y_2=0$) and the gauge transformation that makes its links $1$. A link
$e$ of direction $\mu$ off the tree then equals the holonomy of the closed
path formed by $e$ and the tree paths from the origin to its endpoints; after
cancelling backtracks and inserting one, this path is a product of at most
three conjugated boundaries of rectangles of width $1$ and length at most $L$ in
coordinate planes. By (ii) and $\|XY-1\|\le\|X-1\|+\|Y-1\|$ for unitaries,
$\|V_e-1\|\le3L\max_p\|1-V_p\|$ in this gauge, the maximum over the
plaquettes of the box. With (i), Jensen's inequality and translation
invariance, $\omega_k(\max_p\|1-V_p\|)\le\sum_p\omega_k((2(N-\ReTr V_p))^{1/2})
\le c_L(K_E/s_k^2)^{1/2}\to0$, where $V=\Phi_{s_k}(U)$. Since $O$ is gauge
invariant and uniformly continuous, $|O(V)-O(\mathbf 1)|\le\varpi_O(3L\max_p
\|1-V_p\|)$, which tends to zero in $L^1(\omega_k)$; and $|O_k(f)-O(\mathbf 1)
\sum_xa_k^4f(a_kx)|\le\sum_xa_k^4|f(a_kx)||O_x\circ\Phi_{s_k}-O(\mathbf 1)|$.
\end{proof}
\subsubsection{Which positivity survives}

Theorem~\ref{rF13:tight}(v) is the positivity of a probability measure; it
holds for every flowed limit. Reflection positivity is a different
statement. Let $\theta(x_0,x')=(-x_0,x')$ and let $O\in\mathcal O$ be real
with $O^\theta=O$, so that $\vartheta(O_k(f))=O_k(f\circ\theta)$ for the site
reflection $\vartheta$ by Lemma~\ref{rF13:flow}(ii). Let $S$ be the kernel of
Theorem~\ref{rF13:tight}(iii) for the pair $(O,O)$ in a flowed limit. We say
the limit is \emph{time reflection positive for $O$} if
\begin{equation}\label{rF13:eq:rp}
\iint\overline{f(\theta x)}\,S(x,y)\,f(y)\,dx\,dy\ge0
\quad\text{for all }f\in C_c(\{x_0>0\}) .
\end{equation}
This is the limit of $\omega_k(\vartheta(\overline{O_k(f)})O_k(f))=\nu_k(
\vartheta(\bar F)F)$, $F=\sum_xa_k^4f(a_kx)O_x$ a positive-time function and
$\nu_k=(\Phi_{s_k})_*\omega_k$ the law of the flowed field; so
\eqref{rF13:eq:rp} holds whenever the measures $\nu_k$ are reflection positive
for large $k$, and fails only if they are not. By
Theorem~\ref{rF13:tight}(iv), $S(x,y)=W(x-y)$ almost everywhere for some
$W\in L^\infty(\R^4)$ (write $S(z+y,y)$ and average in $y$ against a
mollifier), and $W(-z)=W(z)$ because $S$ is symmetric. For real $h\in
C_c(\R^3)$ put
\[
C_h(u)=\iint h(x')h(y')\,W(u,x'-y')\,dx'dy' ,
\]
a bounded function of $u\in\R$ defined almost everywhere, with
$C_h(-u)=C_h(u)$.

\begin{lemma}\label{rF13:laplace}
Let $C$ be a bounded continuous function on $(0,\infty)$ with
$\iint\overline{g(u)}g(u')C(u+u')\,du\,du'\ge0$ for all $g\in C_c((0,\infty))$.
Then there is a finite positive measure $\mu$ on $[0,\infty)$ with
$C(u)=\int e^{-uE}\,d\mu(E)$ for all $u>0$.
\end{lemma}

\begin{proof}
On $C_c((0,\infty))$ the form $B(g,g')=\iint\bar g(u)g'(u')C(u+u')\,du\,du'$
is positive semidefinite, hence Hermitian. For $\sigma\ge0$ let $(T_\sigma g)(u)=g(u-
\sigma)$; then $B(T_\sigma g,g')=B(g,T_\sigma g')$, and by the Cauchy--Schwarz
inequality, iterated,
\[
B(T_\sigma g,T_\sigma g)=B(g,T_{2\sigma}g)\le B(g,g)^{1-2^{-n}}
B(T_{2^n\sigma}g,T_{2^n\sigma}g)^{2^{-n}},
\]
with $B(T_\tau g,T_\tau g)\le\|C\|_\infty\|g\|_{L^1}^2$; letting $n\to\infty$,
$T_\sigma$ is a contraction for $B$. On the completion $\mathcal H$ of
$C_c((0,\infty))$ modulo the null space, the $T_\sigma$ are commuting
self-adjoint contractions with $T_\sigma T_\tau=T_{\sigma+\tau}$ and
$T_{2\sigma}=T_\sigma^2\ge0$. By uniqueness of positive roots, $T_\sigma=T_1^\sigma$
for dyadic $\sigma$, so with $\nu_g$ the spectral measure of $T_1$ at $g$
(Lemma~\ref{lem:spectralmeasure}, up to normalisation) $B(g,T_\sigma g)=
\int\lambda^\sigma d\nu_g(\lambda)$. Let $g_\varepsilon\ge0$ be supported in
$[\varepsilon,2\varepsilon]$ with $\int g_\varepsilon=1$, and $\mu_\varepsilon$
the image of $\nu_{g_\varepsilon}$ under $\lambda\mapsto-\log\lambda\in[0,
\infty]$, a measure of mass $B(g_\varepsilon,g_\varepsilon)\le\|C\|_\infty$.
For $\sigma>0$, $B(g_\varepsilon,T_\sigma g_\varepsilon)=\iint g_\varepsilon
g_\varepsilon C(u+u'+\sigma)\to C(\sigma)$ by continuity. A subsequence of
$\mu_\varepsilon$ converges weakly on the compact space $[0,\infty]$
(Lemma~\ref{lem:helly}, after normalisation), and $E\mapsto e^{-\sigma E}$ is
continuous there and vanishes at $\infty$; so $C(\sigma)=\int_{[0,\infty)}
e^{-\sigma E}d\mu(E)$ for dyadic $\sigma>0$, and for all $\sigma>0$ because
both sides are continuous.
\end{proof}

\begin{theorem}\label{rF13:smooth}
Let a flowed limit be time reflection positive for $O$, and let $h\in
C_c(\R^3)$ be real.

\textup{(i)} If $C_h$ agrees almost everywhere on $(0,\infty)$ with a
continuous function, that function is the Laplace transform of a finite
positive measure on $[0,\infty)$.

\textup{(ii)} If $C_h$ agrees almost everywhere on $\R$ with a continuously
differentiable function, then $C_h$ is almost everywhere equal to a constant
$c$. Then, with $m$ the constant one-point kernel of $O$ and real $g\in
C_c(\R)$, the covariance of $X^O(g\otimes h)$ and $X^O(g(\cdot-\sigma)\otimes
h)$ equals $(c-m^2(\int h)^2)(\int g)^2$ for every $\sigma\in\R$: either
$X^O(g\otimes h)$ is almost surely constant for all $g$, or its time
translates are correlated by the same amount at every separation.
\end{theorem}

\begin{proof}
For $f=g\otimes h$ with $g\in C_c((0,\infty))$ the substitution $x=\theta(u,x')$
turns the left side of~\eqref{rF13:eq:rp} into $\iint\bar g(u)g(u')C_h(u+u')
\,du\,du'$, using $W(-z)=W(z)$. Part (i) is Lemma~\ref{rF13:laplace}. For
(ii), let $C$ be the differentiable function; it is even, so $C'(0)=0$, and by
(i) $C(u)=\int e^{-uE}d\mu(E)$ for $u>0$. For $u\ge u_0>0$ the integrand
$Ee^{-uE}$ is bounded by $(eu_0)^{-1}$, so $-C'(u)=\int Ee^{-uE}d\mu(E)$, which
increases to $\int E\,d\mu(E)$ as $u\downarrow0$ by monotone convergence. Since
$C'$ is continuous at $0$, $\int E\,d\mu=0$, so $\mu$ is a multiple of the
point mass at $0$ and $C$ is constant on $(0,\infty)$, hence on $\R$. The
covariance is $\iint g(u)g(u'-\sigma)C_h(u-u')\,du\,du'-(m\int h\int g)^2$.
\end{proof}

The flow is meant to make the flowed field smooth on the scale $\sqrt{8t}$,
and a smooth stationary field has two-point functions that are
differentiable across zero separation. Theorem~\ref{rF13:smooth} says that
such smoothness and time reflection positivity together leave only fields
without decay of time correlations, a property excluded by clustering. The
next proposition shows that this is what happens for the free field, both in
the continuum and on the lattice.

\begin{proposition}\label{rF13:freeflow}
Let $m>0$ and $t>0$, let $\varphi$ be the centred Gaussian field on $\Z^4$
with covariance $(-\Delta+m^2a_k^2)^{-1}$, and let $\nu_k$ be the law of
$e^{s_k\Delta}\varphi$, $s_k=t/a_k^2$, the gradient flow of the Dirichlet
energy at physical time $t$. For $f\in C_c(\R^4)$ put $\varphi_k(f)=a_k^{-1}
\sum_xa_k^4f(a_kx)(e^{s_k\Delta}\varphi)(x)$.

\textup{(i)} $\E\varphi_k(f)\varphi_k(g)\to\iint f(x)W_t(x-y)g(y)$, where
$\hat W_t(p)=e^{-2t|p|^2}(|p|^2+m^2)^{-1}$.

\textup{(ii)} $W_t$ violates~\eqref{rF13:eq:rp}: there is a real $f\in
C_c(\{x_0>0\})$ with $\iint f(\theta x)W_t(x-y)f(y)<0$. Hence the measures
$\nu_k$ are not reflection positive for the site reflection when $k$ is
large, although the unflowed field is; and for $t=0$ the function $C_h$ is
$\int|\hat h(p')|^2e^{-\epsilon|u|}(2\epsilon)^{-1}d^3p'/(2\pi)^3$,
$\epsilon=(|p'|^2+m^2)^{1/2}$, a Laplace transform on $(0,\infty)$.
\end{proposition}

\begin{proof}
(i) In lattice momenta $q\in(-\pi/a_k,\pi/a_k]^4$, with $\hat q^2=\sum_\mu
4a_k^{-2}\sin^2(a_kq_\mu/2)$ and $\hat f_k(q)=\sum_xa_k^4f(a_kx)e^{-ia_kq\cdot x}$,
\[
\E\varphi_k(f)\varphi_k(g)=\int_{(-\pi/a_k,\pi/a_k]^4}\overline{\hat f_k(q)}
\hat g_k(q)\frac{e^{-2t\hat q^2}}{\hat q^2+m^2}\frac{d^4q}{(2\pi)^4},
\]
since $e^{s_k\Delta}$ and $(-\Delta+m^2a_k^2)^{-1}$ act by $e^{-s_ka_k^2\hat q^2}$
and $a_k^{-2}(\hat q^2+m^2)^{-1}$. Riemann sums give $\hat f_k\to\hat f$
pointwise, $|\hat f_k|\le\sum_xa_k^4|f(a_kx)|$ is bounded, and $\sin(y/2)\ge
y/\pi$ on $[0,\pi]$ gives $e^{-2t\hat q^2}\le e^{-8t|q|^2/\pi^2}$; dominated
convergence gives the limit.

(ii) $W_t$ is real, even and, being the Fourier transform of an integrable
function with all moments, smooth with bounded derivatives. For real $h\ne0$,
\[
C_h(u)=\int\frac{d^3p'}{(2\pi)^3}|\hat h(p')|^2\int\frac{dp_0}{2\pi}\,
\frac{e^{ip_0u-2t(p_0^2+|p'|^2)}}{p_0^2+|p'|^2+m^2}
\]
is continuously differentiable, even, positive at $u=0$ and tends to $0$ as
$|u|\to\infty$ by the Riemann--Lebesgue lemma, so it is not constant, and
Theorem~\ref{rF13:smooth}(ii) shows that~\eqref{rF13:eq:rp} fails for some
$f=g\otimes h$. Since $\vartheta$ maps $a_kx$ to $\theta(a_kx)$, the lattice
forms $\E[\varphi_k(f\circ\theta)\varphi_k(f)]$ converge to the left side of
\eqref{rF13:eq:rp} by (i), so they are negative for large $k$, while
$\varphi_k(f)$ is a linear function of the flowed field at positive times.
For $t=0$ the $p_0$-integral is $e^{-\epsilon|u|}/(2\epsilon)$.
\end{proof}

The same computation applies to the field strength of the free Maxwell field,
whose gradient flow is the heat flow $F\mapsto e^{t\Delta}F$ on closed
$2$-forms. What survives is positivity for observables that do not straddle
the plane. On the lattice it is exact for the half-space flow.

\begin{theorem}\label{rF13:half}
Let $\omega_k$ be torus limits as in Lemma~\textup{\ref{rE:torus}} at
arbitrary couplings and $s_k\ge0$ arbitrary. Let $\mathcal D$ be a countable
set of functions in $C_c(\{x_0>0\})$ closed under the dyadic spatial
translations. For $O\in\mathcal O$, $f\in\mathcal D$ and $k$ so large that the
links of every $O_x$ with $f(a_kx)\ne0$ lie in $E_1$, put
\[
O^+_k(f)=\sum_xa_k^4f(a_kx)\bigl(O_x\circ\Phi^{(1)}_{s_k}-\omega_k(O_x\circ
\Phi^{(1)}_{s_k})\bigr),
\]
and $O^+_k(f)=0$ for the finitely many smaller $k$. For loops $C\in\mathcal
C$ in $\{x_0>0\}$ let $W^{(1)}_k(C)$ be the Wilson loop $N^{-1}\Tr V(C/a_k)$
of the half-space flowed field $V=\Phi^{(1)}_{s_k}(U)$, and
$W^+_k(C)=W^{(1)}_k(C)-\omega_k(W^{(1)}_k(C))$, with the same convention.
Let $\mathcal T$ be the vector space with basis the symbols $[O,f,\sigma]$,
$[C,\sigma]$ for dyadic $\sigma>0$, and $[1]$, and let $O_k$ be the linear
map with $O_k[O,f,\sigma]=\tau_{\sigma/a_k}O^+_k(f)$,
$O_k[C,\sigma]=\tau_{\sigma/a_k}W^+_k(C)$ and $O_k[1]=1$.
Definition~\ref{def:limit} and the proof of Theorem~\ref{thm:avgtransfer}
use only the linear structure of the space of test functions, and we apply
them with this $\mathcal T$.

\textup{(i)} $O_k$ takes values in uniform limits of gauge-invariant
continuous functions of finitely many links at positive times, and every
sequence of $k$ has a subsequence along which $(\omega_k,O_k)$ has a
continuum limit in the sense of Definition~\textup{\ref{def:limit}}, with no
renormalisation constants. The dyadic spatial translations act on it by
unitaries commuting with $H$.

\textup{(ii)} The limit is nontrivial if and only if for some basis
element $[O,f,\sigma]$ or $[C,\sigma]$, with $F_k$ the corresponding
$O^+_k(f)$ or $W^+_k(C)$,
\[
\lim_k\omega_k\bigl(\vartheta(\bar F_k)\,\tau_{2\sigma/a_k}F_k\bigr)>0
\]
along the subsequence.

\textup{(iii)} Under the hypothesis of Theorem~\textup{\ref{thm:avgtransfer}},
with $\omega_k$ the torus limits there,
$\sigma(H)\subset\{0\}\cup[m_*,\infty)$ with $m_*=16\varpi/(625e^2\ell)$,
$0$ is a simple eigenvalue, and if the limit is nontrivial then
$\sigma(H)\cap[m_*,\infty)\ne\emptyset$.
\end{theorem}

\begin{proof}
Reflection positivity is exact here for a structural reason: $O\circ
\Phi^{(1)}_{s_k}$ is a function of the links of $E_1$ alone, that is of
links with both endpoints at positive times, and $\omega_k$ is reflection
positive for such functions by Lemma~\ref{rE:torus}(iii).

(i) By Lemma~\ref{rF13:flow}(v), $O^+_k(f)$ and $W^+_k(C)$ are uniform limits of
gauge-invariant continuous functions of finitely many links of $E_1$, whose
endpoints are at times $\ge1$; so the reconstruction of
Appendix~\ref{app:transfer} applies to them (its footnote), with
$\|\Psi_F\|\le\|F\|_{L^2(\omega_k)}\le2\|O\|_\infty\|f\|_{1,k}$ or $\le2$. As in
the proof of Theorem~\ref{thm:flowlimit}, $\Psi_{\tau_nF}=\mathbb T_k^n\Psi_F$, so
the vector of $O_k[O,f,\sigma]$ is $\mathbb T_k^{\sigma/a_k}\Psi_{O^+_k(f)}$.
The numbers $\omega_k(\vartheta(\overline{O_k(b)})\tau_{\rho/a_k}O_k(b'))$ for
basis elements $b,b'$ and dyadic $\rho\ge0$ are countably many and bounded,
so a diagonal subsequence makes them converge. From here the construction
of Theorem~\ref{thm:flowlimit}(iii) applies verbatim with the translates
$\delta_y$ replaced by the basis elements and $T(\rho)$ acting by $\sigma\mapsto
\sigma+\rho$: the limit form is positive semidefinite, the $T(\rho)$ are
commuting positive contractions, each basis element other than $[1]$ is
$T(\sigma/2)$ applied to another one and so is orthogonal to the common kernel,
and $[1]$ spans a vector $\Omega$ fixed by $T(\rho)$ and orthogonal to the
others because the packets are centred. The spatial flow is replaced by
$\Phi^{(1)}_s$, which commutes with spatial translations by
Lemma~\ref{rF13:flow}(ii), and $\mathcal D$ is invariant under them.

(ii) For $b=[O,f,\sigma]$ the squared norm of $\psi(b)$ is the limit of
$\|\mathbb T_k^{\sigma/a_k}\Psi_{O^+_k(f)}\|^2$, which is the displayed one, and $\psi(b)\perp
\Omega$; the limit is nontrivial if and only if one of these vectors is
nonzero.

(iii) The proof of Theorem~\ref{thm:avgtransfer} uses only
Definition~\ref{def:limit}, the linearity of $O_k$ and the gap of $\mathbb
T_k$ on $\Omega_k^\perp$ given by Corollary~\ref{cor:avggap}.
\end{proof}

The half-space flow is not translation covariant in time: $T(\rho)$ moves the
packet and the frozen boundary together, so the vectors $\psi[O,f,\sigma]$ are
states prepared by flowing inside a half-space, not the values of a field.
Theorem~\ref{rF13:half} gives a Hilbert space and a Hamiltonian for every
tuning and without renormalisation; under~\ref{rE:I1}, which is not proved
along any tuning with $\beta_k\to\infty$, the spectral bound in (iii) is
the lattice gap of Corollary~\ref{cor:avggap} read on these vectors, and it
is not a gap of a continuum field theory. Nontriviality is the open
inequality (ii). For the full flow at a
positive physical flow time nothing of this kind is proved. Its
reflection positivity holds up to small errors only when the flow time is
small against the distance to the plane in lattice units.

\begin{proposition}\label{rF13:approx}
Let $\omega_k$ be torus limits as in Lemma~\textup{\ref{rE:torus}}. Let
$O\in\mathcal O$ be real with $O^\theta=O$, supported in $x+[-r,r]^4$
when translated to $x$, and Lipschitz: $|O(V)-O(V')|\le L_O\sum_e\|V_e-V'_e\|$.
Let $f\in C_c(\R^4)$ be supported in $\{x_0\ge d\}$, $d>0$, and $R_k=\lceil d/
a_k\rceil-r-1$. If $\kappa s_k\le R_k/7$, then
\[
\omega_k\bigl(\vartheta(\overline{O_k(f)})O_k(f)\bigr)\ge-2^{4-R_k}(2r+1)^4
\|O\|_\infty L_O\|f\|_{1,k}^2 .
\]
Consequently, if $s_ka_k\to0$, every flowed limit of
Theorem~\textup{\ref{rF13:tight}} along such $s_k$ is time reflection
positive for every such $O$, and for loops $C_1,\dots,C_n\in\mathcal C$ in
$\{x_0>0\}$ the matrix $\bigl(\E[\overline{W(\theta C_i)}W(C_j)]\bigr)_{ij}$
is positive semidefinite, $\theta C_i$ being the reflected loop with the
transported orientation.
\end{proposition}

\begin{proof}
The links of $O_x$ with $f(a_kx)\ne0$ have both endpoints at times $\ge
1+R_k$, and $O$ depends on at most $4(2r+1)^4$ links. By
Lemma~\ref{rF13:flow}(iv), $\sup_U|O_k(f)-O^{(1)}_k(f)|\le\|f\|_{1,k}L_O4(2r+1)^4
2^{1-R_k}$, where $O^{(1)}_k(f)$ is the packet flowed by $\Phi^{(1)}_{s_k}$,
for which $\omega_k(\vartheta(\overline{O^{(1)}_k(f)})O^{(1)}_k(f))\ge0$ by
reflection positivity. For bounded $A,B$, the invariance of $\omega_k$ under
$\vartheta$ gives $|\omega_k(\vartheta(\bar A)A)-\omega_k(\vartheta(\bar B)B)|\le
(\|A\|_\infty+\|B\|_\infty)\sup|A-B|$, with $\|A\|_\infty,\|B\|_\infty\le\|O\|_\infty
\|f\|_{1,k}$; together these give the bound. If $s_ka_k\to0$, then $\kappa
s_k\le R_k/7$ for large $k$ at every $d>0$, and the bound tends to zero. For
the loops, let them lie in $\{x_0\ge d\}$ and put $F=\sum_jc_jW_k(C_j)$ and
$R'_k=\lceil d/a_k\rceil-1$. A Wilson loop is a function of $|C_j|/a_k$
links with Lipschitz constant $1$, because $|\Tr X|\le N\|X\|$ and a product
of unitaries changes in norm by at most the sum of the changes of its
factors; so Lemma~\ref{rF13:flow}(iv) with $h=1$ bounds the change of $F$
under $\Phi_{s_k}\mapsto\Phi^{(1)}_{s_k}$ by $\sum_j|c_j||C_j|a_k^{-1}2^{1-R'_k}
\to0$, while the half-space flowed combination has
$\omega_k(\vartheta(\bar F)F)\ge0$. Since $\vartheta(\overline{W_k(C)})=
\overline{W_k(\theta C)}$, the lattice forms converge to $\sum_{i,j}\bar c_i
c_j\E[\overline{W(\theta C_i)}W(C_j)]$ by Theorem~\ref{rF13:tight}(i).
\end{proof}

So the proved reflection positivity of the four-dimensional flow is confined
to $s_k=o(a_k^{-1})$, physical flow times $t_k=s_ka_k^2=o(a_k)$ and smoothing
radii $\sqrt{t_k}\to0$: it is a statement about the limit of vanishing flow
time, where renormalisation returns.
\subsubsection{Locality, clustering and rotations}

The ballistic bound of Lemma~\ref{rF13:flow}(iv) gives nothing at
$s_k=t/a_k^2$. Let $\underline\Phi{}^{(h)}_s=\vartheta\circ\Phi^{(-h)}_s\circ
\vartheta$, the flow of the links with both endpoints at times $\le h$ by the
plaquettes among them, the others frozen. For $O\in\mathcal O$ and $d>0$ let
$\varepsilon_O(d)$ be the $\limsup_k$ of the supremum of
$\omega_k(|O_x\circ\Phi_{s_k}-O_x\circ\Phi^{(h)}_{s_k}|)$ over all $x\in\Z^4$
and $h\in\Z$ such that every link of $O_x$ has both endpoints at times
$\ge h+d/a_k$, and of $\omega_k(|O_x\circ\Phi_{s_k}-O_x\circ
\underline\Phi{}^{(h)}_{s_k}|)$ over those with every link at times $\le
h-d/a_k$. It is nonincreasing in $d$, at most $2\|O\|_\infty$, and by the
invariance of $\omega_k$ under $\vartheta$ the second supremum for $O$ is
the first one for $O^\theta$.
\begin{enumerate}[label=\textup{(Loc$_t$)},leftmargin=*]
\item \emph{Locality on the physical scale.} $\varepsilon_O(d)\to0$ as
$d\to\infty$, for every $O\in\mathcal O$.
\end{enumerate}
If $s_ka_k\to0$, Lemma~\ref{rF13:flow}(iv) gives $\varepsilon_O(d)=0$ for
every $d>0$. At $s_k=t/a_k^2$ the heat kernel suggests (Loc$_t$) with
$\varepsilon_O(d)$ decaying like $e^{-d^2/(8t)}$ up to powers, as the heat
kernel does; no bound uniform in $k$ is known, since it would have to hold
for typical rather than all configurations.

\begin{proposition}\label{rF13:cluster}
Let $\omega_k$ be torus limits as in Lemma~\textup{\ref{rE:torus}} at
couplings $\beta_k\ge0$, and consider a flowed limit of
Theorem~\textup{\ref{rF13:tight}}. Let $O,O'\in\mathcal O$ and $f,g\in
C_c(\R^4)$.

\textup{(i)} If $O$ is real with $O^\theta=O$ and $f$ is real and supported
in $\{x_0\ge d\}$, $d>0$, then
$\iint f(\theta x)S(x,y)f(y)\,dx\,dy\ge-2\|O\|_\infty\|f\|_{L^1}^2
\varepsilon_O(d/2)$.

\textup{(ii)} Under \textup{\ref{rE:I1}}, if $g$ is supported in $\{x_0\le0\}$
and $f$ in $\{x_0\ge r\}$, $r>0$, then for every $d\in(0,r/2)$
\[
\begin{aligned}
&\frac{\bigl|\Cov\bigl(X^{O'}(g),X^O(f)\bigr)\bigr|}{\|f\|_{L^1}\|g\|_{L^1}}
\le\|O\|_\infty\|O'\|_\infty e^{-m_*(r-2d)}\\
&\qquad+2\|O\|_\infty\varepsilon_{O'}(d)+2\|O'\|_\infty\varepsilon_O(d).
\end{aligned}
\]
Under \textup{(Loc$_t$)} the right side tends to $0$ as $r\to\infty$; if
$\varepsilon_O(d),\varepsilon_{O'}(d)\le c\,e^{-d/\xi}$, it is at most
$C e^{-m_*r/(1+2m_*\xi)}$.

\textup{(iii)} Let $f\in\mathcal D$ be supported in $\{x_0\ge d\}$ and
$\sigma>0$ dyadic, let $\psi[O,f,\sigma]$ be the vector of
Theorem~\textup{\ref{rF13:half}} along a common subsequence, and
$F_k=O_k(f)-\omega_k(O_k(f))$. Then the limit $\Lambda=\lim_k\omega_k(\vartheta(\bar F_k)\tau_{2\sigma/a_k}F_k)$
satisfies
\[
\bigl|\Lambda-\|\psi[O,f,\sigma]\|^2\bigr|\le8\|O\|_\infty\|f\|_{L^1}^2\,
\varepsilon_O(d/2),
\]
and $\Lambda$ is a covariance of the flowed limit, with
$\Cov(X,Y)=\E[XY]-\E X\,\E Y$:
\[
\Lambda=\Cov\bigl(\overline{X^{O^\theta}(f\circ\theta)},X^O(f(\cdot-2\sigma
e_0))\bigr).
\]
\end{proposition}

\begin{proof}
Write $\|f\|_{1,k}=\sum_xa_k^4|f(a_kx)|$, which tends to $\|f\|_{L^1}$, and
let $r_O$ be such that the links of $O_x$ lie within distance $r_O$ of $x$.
For bounded $A,B$ and a state invariant under $\vartheta$,
\begin{equation}\label{rF13:eq:swap}
\begin{aligned}
&|\omega(\vartheta(\bar A)A')-\omega(\vartheta(\bar B)B')|\\
&\qquad\le\|A'\|_\infty\,\omega(|A-B|)+\|B\|_\infty\,\omega(|A'-B'|),
\end{aligned}
\end{equation}
because $\omega(|\vartheta(X)|)=\omega(|X|)$.

(i) Let $A=O_k(f)$ and let $B$ be the same packet flowed by $\Phi^{(1)}_{s_k}$.
The links of $O_x$ with $f(a_kx)\ne0$ are at times $\ge d/a_k-r_O$, which
exceeds $1+d/(2a_k)$ for large $k$, so $\omega_k(|A-B|)\le\|f\|_{1,k}\delta_k$
with $\limsup_k\delta_k\le\varepsilon_O(d/2)$. The function $B$ depends only
on links of $E_1$, so $\omega_k(\vartheta(\bar B)B)\ge0$ by reflection
positivity of the torus limit, Lemma~\ref{rE:torus}(iii), which passes to
uniform limits by Lemma~\ref{rF13:flow}(v). By
\eqref{rF13:eq:swap} with $A'=A$, $B'=B$, and $\|A\|_\infty,\|B\|_\infty\le
\|O\|_\infty\|f\|_{1,k}$, we get $\omega_k(\vartheta(\bar A)A)\ge-2\|O\|_\infty
\|f\|_{1,k}^2\delta_k$. Since $\vartheta(\bar A)=O_k(f\circ\theta)$ by
Lemma~\ref{rF13:flow}(ii), the left side converges to the integral by
Theorem~\ref{rF13:tight}(i),(iii).

(ii) Let $A=O'_k(g)$, $B=O_k(f)$, $h_A=r_{O'}+\lceil d/a_k\rceil$ and $h_B=
\lceil r/a_k\rceil-r_O-\lceil d/a_k\rceil$. Let $A'$ be the packet $A$ with
$\Phi_{s_k}$ replaced by $\underline\Phi{}^{(h_A)}_{s_k}$ and $B'$ the packet
$B$ with $\Phi_{s_k}$ replaced by $\Phi^{(h_B)}_{s_k}$. The links of $O'_x$
with $g(a_kx)\ne0$ lie at times $\le r_{O'}\le h_A-d/a_k$, and those of
$O_x$ with $f(a_kx)\ne0$ at times $\ge h_B+d/a_k$; so $\omega_k(|A-A'|)\le
\|g\|_{1,k}\delta'_k$ and $\omega_k(|B-B'|)\le\|f\|_{1,k}\delta_k$ with
$\limsup\delta'_k\le\varepsilon_{O'}(d)$, $\limsup\delta_k\le\varepsilon_O(d)$.
For bounded functions, $|\Cov(A,B)-\Cov(A',B')|\le2\|B\|_\infty\omega(|A-A'|)
+2\|A'\|_\infty\omega(|B-B'|)$. Now $A'$ depends on links with both endpoints
at times $\le h_A$ and $B'$ on links at times $\ge h_B$. After the time
translation by $-h_A-1$, which preserves $\omega_k$, $A'$ becomes
$\vartheta(\bar F)$ with $F$ a function of the links of $E_1$, and $B'$
becomes $\tau_nG$ with $G$ a function of the links of $E_1$ and
$n=h_B-h_A-2$. With the transfer matrix $\mathbb T_k$ of
Lemma~\ref{rE:torus}(iv), $\Cov(A',B')=\ip{\Psi_F}{\mathbb T_k^n(\Psi_G-
\ip\Omega{\Psi_G}\Omega)}$, and by Corollary~\ref{cor:avggap} under
\ref{rE:I1} the spectrum of $\mathbb T_k$ on $\Omega^\perp$ lies in
$[0,e^{-m_*a_k}]$, as in the proof of Theorem~\ref{thm:avgtransfer}. Hence
$|\Cov(A',B')|\le e^{-m_*a_kn}\|F\|_{L^2}\|G\|_{L^2}\le e^{-m_*a_kn}
\|O\|_\infty\|O'\|_\infty\|f\|_{1,k}\|g\|_{1,k}$, and $a_kn\ge r-2d-a_k(r_O+
r_{O'}+4)$. The covariances converge along the subsequence by
Theorem~\ref{rF13:tight}(i). The last claims follow with $d=r/3$, and with
$d=m_*\xi r/(1+2m_*\xi)$, for which $m_*(r-2d)=d/\xi=m_*r/(1+2m_*\xi)$.

(iii) Let $B_k=O^+_k(f)$ be the centred half-space packet, so that
$\|\psi[O,f,\sigma]\|^2=\lim_k\omega_k(\vartheta(\bar B_k)\tau_{2\sigma/a_k}B_k)$
(proof of Theorem~\ref{rF13:half}(ii)). As in (i), $\omega_k(|F_k-B_k|)\le
2\|f\|_{1,k}\delta_k$ with $\limsup\delta_k\le\varepsilon_O(d/2)$, the centring
doubling the error, and $\|F_k\|_\infty,\|B_k\|_\infty\le2\|O\|_\infty
\|f\|_{1,k}$. Apply \eqref{rF13:eq:swap} with $A=F_k$, $A'=\tau_{2\sigma/
a_k}F_k$, $B=B_k$, $B'=\tau_{2\sigma/a_k}B_k$, using translation invariance.
The identification of the limit is Lemma~\ref{rF13:flow}(ii) and
Theorem~\ref{rF13:tight}(i).
\end{proof}

Part (iii) gives the inequality of Theorem~\ref{rF13:half}(ii) from a lower
bound on time correlations of the full flowed limit: if for some $f$
supported in $\{x_0\ge d\}$ the reflected covariance exceeds
$8\|O\|_\infty\|f\|_{L^1}^2\varepsilon_O(d/2)$, the half-space limit is
nontrivial. Both the lower bound and (Loc$_t$) are open.

We turn to rotations. Let $\mathcal O_0\subset\mathcal O$ be the observables
with $O^g=O$ for all $g\in B_4$; the gauge-invariant sums over the
$B_4$-orbit of a local function belong to it. For a flowed limit let
$G_S$ be the set of $R\in\mathrm O(4)$ such that for all $n$, $O_j\in
\mathcal O_0$ and $f_j\in C_c(\R^4)$ the family $(X^{O_j}(f_j\circ
R^{-1}))_{j\le n}$ has the law of $(X^{O_j}(f_j))_{j\le n}$, and the same with
$\mathfrak E$ among the $O_j$ when the limit is taken under (E$_t$) as in
Theorem~\ref{rF13:tight}(vi). Let $R_\theta$ be the rotation by $\theta$ in
the $(x_0,x_1)$-plane with $\cos\theta=\frac35$, $\sin\theta=\frac45$.

\begin{theorem}\label{rF13:o4}
\textup{(i)} $G_S$ is a closed subgroup of $\mathrm O(4)$ containing $B_4$.
It is either finite or equal to $\mathrm O(4)$, and it equals
$\mathrm O(4)$ if and only if it contains an element of infinite order, if
and only if $R_\theta\in G_S$. The group generated by $B_4$ and $R_\theta$
is dense in $\mathrm O(4)$, and its rotations are dense in $\SO(4)$.

\textup{(ii)} $R\in G_S$ if and only if, along the subsequence, for all
$n$, $O_j$, $f_j$ as above and every bounded Lipschitz $\varphi$ on $\C^n$,
\begin{equation}\label{rF13:eq:Wt}
\lim_k\Bigl[\omega_k\bigl(\varphi((O_{j,k}(f_j\circ R^{-1}))_j)\bigr)
-\omega_k\bigl(\varphi((O_{j,k}(f_j))_j)\bigr)\Bigr]=0 .
\end{equation}
Let \textup{(W$_t$)} be the statement that~\eqref{rF13:eq:Wt} holds for
$R=R_\theta$ along the full sequence, for all such data. Then
\textup{(W$_t$)} holds if and only if $G_S=\mathrm O(4)$ for every flowed
limit along a subsequence.

\textup{(iii)} Let $s_k=t/a_k^2$ and assume \textup{(E$_t$)}. If $C,C'\in
\mathcal C$ and $C'$ arises from $C$ by replacing a segment of $C$ that is
one side, or two consecutive sides, of a rectangle $Q$ in a coordinate
plane by the remaining sides of $Q$, then
\[
\E|W(C)-W(C')|\le\sqrt{K_E/6}\;|Q|/t ,
\]
$|Q|$ the area. Consequently the Wilson loops of rectangles with dyadic
corners in the coordinate planes extend, as limits in $L^1$, to rectangles
with arbitrary real corners, with $\E|W(\partial Q)-W(\partial(Q+v))|\le
\sqrt{K_E/6}\,|v|_1P(Q)/t$, $P$ the perimeter, and their joint laws with the
$X^O(f)$ are invariant under all translations of $\R^4$.
\end{theorem}

\begin{proof}
(i) If $R,R'\in G_S$, applying the definition to $f_j\circ R'^{-1}$ and then
to $f_j$ shows $R'R\in G_S$, and applying it to $f_j\circ R$ shows $R^{-1}\in
G_S$. If $R_m\to R$ with $R_m\in G_S$, then $f\circ R_m^{-1}\to f\circ R^{-1}$
in $L^1$; by Theorem~\ref{rF13:tight}(ii) $|X^O(f\circ R_m^{-1})-X^O(f\circ
R^{-1})|\le\|O\|_\infty\|f\circ R_m^{-1}-f\circ R^{-1}\|_{L^1}$, and
$\E|X^{\mathfrak E}(f\circ R_m^{-1})-X^{\mathfrak E}(f\circ R^{-1})|\le K_E
\|f\circ R_m^{-1}-f\circ R^{-1}\|_{L^1}$ by Theorem~\ref{rF13:tight}(vi). So
the laws converge and $R\in G_S$. For $g\in B_4$ and $O\in\mathcal O_0$,
Theorem~\ref{rF13:tight}(iv) gives that $(X^O(f\circ g^{-1}))$ has the law of
$(X^O(f))$, and (vi) the same for $\mathfrak E$; so $B_4\subset G_S$.

The next step is the argument of Proposition~\ref{prop:rotations}(i),(iv),
which uses only that the group is closed in $\mathrm O(4)$ and contains
$B_4$. Let $H$ be such a group. It is a compact Lie group, and its Lie
algebra $\mathfrak h\subset\mathfrak{so}(4)$ is invariant under
$\mathrm{Ad}(g)$, $g\in B_4$. With $E_{ij}=e_ie_j^{\mathsf T}-e_je_i^{\mathsf
T}$ and the sign changes $D_\epsilon\in B_4$, $\epsilon\in\{\pm1\}^4$,
$\mathrm{Ad}(D_\epsilon)E_{ij}=\epsilon_i\epsilon_jE_{ij}$; the six
characters $\epsilon\mapsto\epsilon_i\epsilon_j$ are distinct, so the
averages $2^{-4}\sum_\epsilon\epsilon_i\epsilon_j\mathrm{Ad}(D_\epsilon)$
project $\mathfrak h$ onto $\mathfrak h\cap\R E_{ij}$ and $\mathfrak h$ is
spanned by the $E_{ij}$ it contains. The permutations of the four axes act
transitively on the pairs $\{i,j\}$, so $\mathfrak h$ is $0$ or
$\mathfrak{so}(4)$. The odd sign changes are needed here: under the
rotations of $B_4$ alone, $\mathfrak{so}(4)$ splits into its self-dual and
anti-self-dual ideals, since $E_{01}$ and $E_{23}$ carry the same character
of the even sign changes. If $\mathfrak h=0$, $H$ is discrete and compact,
hence finite. If $H$ contains an element of infinite order, the closure of
its powers is an infinite compact subgroup, so $H$ is infinite,
$\mathfrak h=\mathfrak{so}(4)$, $H$ contains the identity component
$\SO(4)$, and with the reflections in $B_4$ it is $\mathrm O(4)$. Applied to
$H=G_S$ this gives the first two sentences of (i), once $R_\theta$ is known
to have infinite order: $e^{i\theta}=(3+4i)/5=(2+i)/(2-i)$, and
$e^{in\theta}=1$ with $n\ge1$ would give $(2+i)^n=(2-i)^n$ in $\Z[i]$,
while $2+i$ and $2-i$ are primes of $\Z[i]$ that are not associate (the
associates of $2+i$ are $\pm(2+i)$, $\pm(1-2i)$); this contradicts unique
factorisation. Applied to $H$ the closure of the group $\Gamma$ generated by
$B_4$ and $R_\theta$ it gives $H=\mathrm O(4)$, so $\Gamma$ is dense. The
rotations $\Gamma\cap\SO(4)$ have index at most $2$ in $\Gamma$, so their
closure $K$ satisfies $\mathrm O(4)=K\cup gK$ for some $g$; a closed
subgroup of finite index is open and contains the identity component, so
$K\supseteq\SO(4)$.

(ii) By Theorem~\ref{rF13:tight}(i),(vi), applied to the functions
$f_j\circ R^{-1}\in C_c(\R^4)$ and to $f_j$, both expectations converge to the
expectations of $\varphi$ under the two limit laws, and bounded Lipschitz
functions determine laws. If (W$_t$) holds, $R_\theta$ lies in every $G_S$,
which is then $\mathrm O(4)$ by (i). Conversely, if~\eqref{rF13:eq:Wt}
fails for $R_\theta$ and some data, the difference stays above some
$\varepsilon>0$ in modulus along a subsequence, a further subsequence
carries a flowed limit by Theorem~\ref{rF13:tight}, and for it
$R_\theta\notin G_S$.

(iii) For $k$ larger than the levels, $C/a_k$, $C'/a_k$ are lattice loops and
$Q/a_k$ is a rectangle of $|Q|/a_k^2$ plaquettes. Add its plaquettes one at
a time, row by row starting from the corner at which the replaced segment
begins; each new plaquette shares with the current path a run of one, two or
three consecutive sides, and replacing the run by the other sides of the
plaquette is the insertion of its boundary at the first vertex of the run
followed by cancellation of backtracks. By Lemma~\ref{rF13:stokes}(ii) and the triangle inequality,
$\|V(C/a_k)-V(C'/a_k)\|\le\sum_{p\subset Q/a_k}\|1-V_p\|$ with $V=
\Phi_{s_k}(U)$, both holonomies based at the same point; and $|W_k(C)-
W_k(C')|\le\|V(C/a_k)-V(C'/a_k)\|$ because $|\Tr X|\le N\|X\|$. By
translation and $B_4$ invariance of $\omega_k$ and Lemma~\ref{rF13:flow}(ii),
all plaquettes have the same mean action, which is $\omega_k(e_0\circ
\Phi_{s_k})/6=\omega_k(\mathfrak E_k(0))/(12s_k^2)\le K_E/(12s_k^2)$. By
Lemma~\ref{rF13:stokes}(i) and Jensen's inequality $\omega_k(\|1-V_p\|)\le
(K_E/(6s_k^2))^{1/2}$, and the sum over the plaquettes is at most
$\sqrt{K_E/6}\,|Q|/(a_k^2s_k)=\sqrt{K_E/6}\,|Q|/t$. The function $|w-w'|$ is
bounded and continuous on the product of two discs, so the bound passes to
the limit by Theorem~\ref{rF13:tight}(i). The lattice inequality
$\omega_k(|W_k(C)-W_k(C')|)\le\sqrt{K_E/6}\,|Q|/t$ holds for all closed lattice
paths related in this way, so moves can be chained on the lattice, through
paths with backtracks, before the limit is taken; only the end loops need
to lie in $\mathcal C$.

For rectangles $Q''\subset Q'$ in one coordinate plane with dyadic corners,
$\partial Q'$ arises from $\partial Q''$ by moving the four sides outwards,
four moves of the first kind of total area $|Q'\setminus Q''|$. For two
dyadic rectangles close to a given real rectangle $Q$, compare both with
their intersection: the dyadic approximations of $Q$ are Cauchy in $L^1$,
and the limit $W(\partial Q)$ does not depend on them. For $v=ue_\rho$ with
$e_\rho$ orthogonal to the plane of $Q$, the loop $\partial(Q+v)$, joined to
$\partial Q$ by the segment from a corner $q$ to $q+v$ and back, arises from
$\partial Q$ by four moves, one across each lateral face $\sigma\times[0,u]
e_\rho$ of the box, of total area $|u|P(Q)$, the connecting segments of
consecutive faces cancelling as backtracks; a translation within the plane
by $ue_\mu$ is a move and the inverse of a move, of total area at most
$|u|P(Q)$; the bound is symmetric in $C,C'$. The trace does not
see the connecting segments. This gives the bound for dyadic $Q$ and $v$,
hence for real ones by the limits, and translation invariance of the joint
laws follows from Theorem~\ref{rF13:tight}(iv) by approximating $v$ by dyadic
vectors.
\end{proof}

\begin{remark}\label{rF13:o4nogo}
Neither the lattice structure nor the flow gives (W$_t$). The flow of a
Gaussian field multiplies its covariance in momentum space by
$e^{-2t|p|^2}$, which is $\mathrm O(4)$-invariant and nowhere zero, so the
flowed field has exactly the rotation symmetry of the unflowed one. The
centred Gaussian field on $\R^4$ with covariance $W(x)=\prod_\mu e^{-m|x_\mu|}/
(2m)$, the product of four Ornstein--Uhlenbeck kernels, is invariant under
translations and $B_4$, its covariance satisfies~\eqref{rF13:eq:rp} in every
coordinate direction (for real $f$ supported in $\{x_0>0\}$ the reflected
form is $\iint e^{-mx_0}f(x)e^{-my_0}f(y)K(x'-y')\,dx\,dy\ge0$ with
$K=\prod_{\mu\ge1}e^{-m|x_\mu|}/(2m)$ positive definite), and it clusters
at rate $m$; but $W$ is not radial, and neither is the covariance of its
flow, $\hat W_t(p)=e^{-2t|p|^2}\prod_\mu(p_\mu^2+m^2)^{-1}$, which by the
argument of Proposition~\ref{rF13:freeflow}(ii) violates~\eqref{rF13:eq:rp}.
Part (iii) of Theorem~\ref{rF13:o4} does
not help either: a rectangle rotated by $R_\theta$ is not rectilinear, and
continuity in the loop says nothing about the law of rotated loops. So
(W$_t$) is a dynamical property of the Wilson measure along the tuning, as
(W$_{\mathrm s}$) is in \ref{rE:I4}, and it is not proved.
\end{remark}
\subsubsection{The area law and Gaussian limits}

Let $s_k=t/a_k^2$ and let $C_{R,T}\in\mathcal C$ be the boundary of
$[0,R]\times[0,T]$ in the $(x_0,x_1)$-plane, $R,T$ dyadic. Consider:
\begin{enumerate}[label=\textup{(A$_t$)},leftmargin=*]
\item \emph{Area law for flowed loops.} There are $\sigma>0$, $c\ge0$ and
$C_0$ with $|\E W(C_{R,R})|\le C_0e^{-\sigma R^2+cR}$ for all dyadic $R\ge1$.
\end{enumerate}
Here $\sigma$ is a string tension in physical units at flow time $t$, and
no renormalisation is involved: $|W|\le1$. The comparison class is the
following. A flowed limit has \emph{abelian Gaussian loops with line
constant $K$} if there are measurable real functions $D^{(j)}_{\mu\nu}$ on
$\R^4$, $1\le j\le N$, with
\[
K=\max_{j,\mu,\nu}\ \sup_{x\in\R^4,\ L}\int_L|D^{(j)}_{\mu\nu}(x-y)|\,dy<
\infty ,
\]
the supremum over the lines $L$ parallel to a coordinate axis, $dy$ the
arc length, and, for all dyadic $R,T$, centred Gaussian random variables
$\alpha_j$ with $\Var\alpha_j=\oint_{C_{R,T}}\oint_{C_{R,T}}D^{(j)}_{\mu\nu}
(x-y)\,dx^\mu dy^\nu$ such that $\E W(C_{R,T})=N^{-1}\sum_j\E e^{i\alpha_j}$.
This holds with $K=0$ for the flat limit $W\equiv1$, and for the Wilson loops
$N^{-1}\Tr\exp\oint_CA$ of every centred Gaussian connection $A=i\,
\mathrm{diag}(A^1,\dots,A^N)$ with values in a Cartan subalgebra whose
components have covariance kernels $D^{(j)}_{\mu\nu}$ with finite line
integrals: the matrices commute, so path ordering plays no role and
$\alpha_j=\oint_CA^j$.

For a closed polygon $C$ put $J_\nu(p)=\oint_Ce^{-ip\cdot y}dy^\nu$, so that
$|J_\nu|\le|C|$ and $p_\nu J_\nu(p)=i\oint_Cd(e^{-ip\cdot y})=0$. A centred
Gaussian $1$-form field with covariance $\hat D_{\mu\nu}\in L^1(\R^4)$ in
momentum space has, by Fubini's theorem, loop variances
\begin{equation}\label{rF13:eq:loopvar}
\begin{aligned}
\Var\oint_CA&=\int\overline{J_\mu(p)}\hat D_{\mu\nu}(p)J_\nu(p)\frac{d^4p}{
(2\pi)^4}\\&=\oint_C\oint_CD_{\mu\nu}(x-y)\,dx^\mu dy^\nu ,
\end{aligned}
\end{equation}
and a term $p_\mu v_\nu(p)+v_\mu(p)p_\nu$ of $\hat D$ does not contribute.
So the loop variances do not depend on the gauge. The \emph{heat-flowed
Maxwell field} at coupling $g$ is the field with $\hat D_{\mu\nu}=g^2
\delta_{\mu\nu}e^{-2t|p|^2}|p|^{-2}$, the Feynman-gauge propagator multiplied
by the flow factor $e^{-2t|p|^2}$ of the linearised flow~\cite[Sec.~2]{Lus10};
the \emph{heat-flowed Proca field} of mass $m>0$ has $\hat D_{\mu\nu}=
e^{-2t|p|^2}(\delta_{\mu\nu}+p_\mu p_\nu/m^2)(|p|^2+m^2)^{-1}$.

\begin{theorem}\label{rF13:area}
\textup{(i)} If a flowed limit has abelian Gaussian loops with line constant
$K$, then $\E W(C_{R,T})\ge e^{-4K(R+T)}$; in particular $\E W(C_{R,R})\ge e^{-8KR}$.

\textup{(ii)} A flowed limit satisfying \textup{(A$_t$)} does not have
abelian Gaussian loops with any finite line constant. In particular
$W(C_{R,R})$ is not almost surely $1$ for large $R$, and the loops are not
those of a heat-flowed Maxwell field in a Cartan subalgebra (components
$A^j$ that are linear combinations of independent heat-flowed Maxwell
fields), at any coupling.

\textup{(iii)} The heat-flowed Maxwell field has $D_{\mu\nu}=g^2\delta_{\mu\nu}
D_t$ with
\[
D_t(z)=\frac{1-e^{-|z|^2/(8t)}}{4\pi^2|z|^2},\qquad K\le\frac{g^2}{2\sqrt2\,
\pi^2\sqrt t} .
\]
The heat-flowed Proca field has the loop variances of $\delta_{\mu\nu}D^P_t$
with $0<D^P_t\le D_t$, hence a perimeter law with the same constant at
$g=1$.
\end{theorem}

\begin{proof}
(i) $\E e^{i\alpha_j}=e^{-\Var\alpha_j/2}$. The loop $C_{R,T}$ is the union of
four sides, each contained in a coordinate line, so for every $x$ and all
$\mu,\nu$, $\oint_{C_{R,T}}|D^{(j)}_{\mu\nu}(x-y)||dy|\le4K$, and since
$|dx^\mu|,|dy^\nu|$ are at most the arc lengths,
$\Var\alpha_j\le\oint\oint|D^{(j)}_{\mu\nu}(x-y)||dx||dy|\le4K|C_{R,T}|=8K(R+T)$.
Hence $\E W\ge N^{-1}\sum_je^{-4K(R+T)}$.

(ii) By (i) and (A$_t$), $e^{-8KR}\le C_0e^{-\sigma R^2+cR}$ for all dyadic
$R\ge1$, which fails for large $R$ because $\sigma>0$. The flat limit has
$K=0$, and the Maxwell field in a Cartan subalgebra has $D^{(j)}$ a
nonnegative multiple of $\delta_{\mu\nu}D_t$, with finite $K$ by (iii).

(iii) Since $e^{-2t|p|^2}|p|^{-2}=\int_{2t}^\infty e^{-\tau|p|^2}d\tau$, the
inverse Fourier transform of $e^{-\tau|p|^2}$ on $\R^4$ is $(4\pi\tau)^{-2}
e^{-|z|^2/(4\tau)}$, and all integrands are positive, Tonelli's theorem gives
\[
D_t(z)=\int_{2t}^\infty\frac{e^{-|z|^2/(4\tau)}}{16\pi^2\tau^2}\,d\tau
=\frac1{4\pi^2|z|^2}\int_0^{|z|^2/(8t)}e^{-u}\,du ,
\]
with $u=|z|^2/(4\tau)$. From $1-e^{-v}\le\min(1,v)$, $D_t(z)\le\min(A,B|z|^{-2})$
with $A=(32\pi^2t)^{-1}$, $B=(4\pi^2)^{-1}$. On a line $L$, with $s$ the
coordinate along $L$ from the foot of the perpendicular from $x$, $|x-y|\ge
|s|$, so $\int_LD_t(x-y)\,dy\le\int_\R\min(A,Bs^{-2})\,ds=4\sqrt{AB}=(2\sqrt2
\pi^2\sqrt t)^{-1}$. For the Proca field, \eqref{rF13:eq:loopvar} and
$p_\nu J_\nu=0$ remove the term $p_\mu p_\nu$, which is integrable against the
other factors, and $e^{-2t|p|^2}(|p|^2+m^2)^{-1}=\int_{2t}^\infty e^{-(\tau-2t)
m^2}e^{-\tau|p|^2}d\tau$ gives $D^P_t(z)=\int_{2t}^\infty e^{-(\tau-2t)m^2}(4\pi
\tau)^{-2}e^{-|z|^2/(4\tau)}d\tau\in(0,D_t(z)]$.
\end{proof}

The free Proca field is a Gaussian field whose covariance decays at the
rate $m$, and its lattice version satisfies the cube condition
(Theorem~\ref{thm:ffproca}); its heat-flowed loops obey a perimeter law by
(iii). So, for Gaussian fields, exponential decay at a positive rate and the
cube condition of the lattice field are compatible with a perimeter law at
positive flow time. These fields are not lattice gauge theories, and the
flow of a Gaussian field is the heat flow, so the example shows only that
(A$_t$) is an input independent of the gap. For the Wilson measure at fixed small $\beta$ the
unflowed Wilson loops obey an area law~\cite{OSei78}, but that is not a
continuum tuning, and no relation between an unflowed string tension along a
tuning and (A$_t$) is proved: at $s_k=t/a_k^2$ the flow is local only
ballistically, Lemma~\ref{rF13:flow}(iv). The unflowed comparison needs
a renormalisation of the perimeter term, since the line integrals of the
unflowed propagator $(4\pi^2|z|^2)^{-1}$ diverge; at positive flow time
$K$ is finite and (A$_t$) is a statement about bounded observables.

\subsubsection{The limit of vanishing flow time}\label{rF13:zero}

One fact relates different flow times.

\begin{lemma}\label{rF13:mono}
Write $\mathfrak E^{t}_k$ for $\mathfrak E_k$ at $s_k=t/a_k^2$. For $0<t\le
t'$ and every $k$, $\omega_k(\mathfrak E^{t'}_k(0))\le(t'/t)^2\omega_k(
\mathfrak E^{t}_k(0))$. So \textup{(E$_t$)} implies \textup{(E$_{t'}$)} with
$K_E(t')\le(t'/t)^2K_E(t)$. Under \textup{(E$_{t_0}$)} one subsequence carries
flowed limits for all dyadic $t\ge t_0$ at once, and along a further
subsequence $\hat E(t)=\lim_kt^{-2}\omega_k(\mathfrak E^t_k(0))$ exists and is
nonincreasing on the dyadic $t\ge t_0$.
\end{lemma}

\begin{proof}
By Lemma~\ref{rF13:flow}(vi), $\omega_k(e_0\circ\Phi_{s'})\le\omega_k(e_0\circ
\Phi_s)$ for $s\le s'$, and $\mathfrak E^t_k(0)=2(t/a_k^2)^2e_0\circ
\Phi_{t/a_k^2}$. The families of Theorem~\ref{rF13:tight} for countably many
flow times form one countable family, so the proof of that theorem applies
to their union; the numbers $t^{-2}\omega_k(\mathfrak E^t_k(0))$ are bounded
for each $t$ and nonincreasing in $t$ at each $k$, and a diagonal
subsequence makes them converge.
\end{proof}

The monotonicity goes the wrong way for $t\to0$: (E$_t$) for small $t$ is the
stronger statement, and in perturbation theory $t^2E(t)$ tends to $0$ like
the running coupling at the scale $\sqrt{8t}$~\cite[(2.32)]{Lus10}. Nothing
proved here controls a flowed limit as $t\to0$. The relation to the
Osterwalder--Schrader framework is as follows.

\begin{enumerate}[label=\textup{(\arabic*)},leftmargin=*]
\item At fixed $t>0$ the full flowed limit is a random field with bounded
correlation kernels, invariant under translations and $B_4$
(Theorem~\ref{rF13:tight}). Reflection positivity, on which the
reconstruction of a Hilbert space and the gap of
Theorem~\ref{thm:avgtransfer} rest, is not proved for it, and it fails for
every limit whose smeared two-point functions are differentiable and decay
(Theorem~\ref{rF13:smooth}), as for the flowed free field. Under (Loc$_t$)
it holds up to $2\|O\|_\infty\|f\|_{L^1}^2\varepsilon_O(d/2)$, and under
(Loc$_t$) and \ref{rE:I1} the limit clusters in time
(Proposition~\ref{rF13:cluster}).
\item The reflection positive objects at fixed $t>0$ are the half-space
flows: their limits exist for every tuning without renormalisation, and
under \ref{rE:I1} they have $\sigma(H)\subset\{0\}\cup[m_*,\infty)$; they are
nontrivial if and only if the inequality of Theorem~\ref{rF13:half}(ii)
holds, which by Proposition~\ref{rF13:cluster}(iii) follows from (Loc$_t$)
and a lower bound on time correlations of the full flow. They are states,
not a field.
\item Exact reflection positivity of the full flow is proved only for
physical flow times $t_k=o(a_k)$ (Proposition~\ref{rF13:approx}). There the
smoothing radius $\sqrt{t_k}$ tends to zero, the flowed fields approach the
unflowed ones, and a nontrivial limit needs the renormalisation constants
and inputs \ref{rE:I2}--\ref{rE:I5} of Sec.~\ref{ssec:rEmissing}, the setting
of Theorem~\ref{rE:suffices}.
\item Passing from $t>0$ to $t=0$, that is, interchanging $a_k\to0$ with
$t\to0$ for renormalised observables and recovering reflection positivity,
$\mathrm O(4)$ invariance and nontriviality of the unflowed limit, is the
content of statement (E); with (D), or with~\ref{rE:I1} along the tuning
and Theorem~\ref{thm:avgtransfer}, it would give the continuum mass gap.
None of this is proved here, and no continuum mass gap is claimed.
\end{enumerate}

The open inputs of this subsection are: (E$_t$), the bounded flowed action
along the tuning; (Loc$_t$), locality of the flow on the physical scale for
typical configurations; (W$_t$), one rotation of infinite order as an
asymptotic symmetry; (A$_t$), the area law for flowed loops; the
nontriviality inequality of Theorem~\ref{rF13:half}(ii); \ref{rE:I1}; and the
limit $t\to0$. None of them is proved here for any tuning with
$\beta_k\to\infty$.
